% Generated draft: requires mathematical and source review.
\documentclass[11pt]{article}
\usepackage[margin=25mm]{geometry}
\usepackage{fontspec}
\setmainfont{Latin Modern Roman}
\usepackage{amsmath,amssymb,mathtools,unicode-math}
\setmathfont{Latin Modern Math}
\usepackage{enumitem,needspace,xurl,hyperref}
\setlength{\parindent}{0pt}
\setlength{\parskip}{5pt}
\setlength{\emergencystretch}{3em}
\newcommand{\N}{\mathbb{N}}
\newcommand{\Z}{\mathbb{Z}}
\newcommand{\Q}{\mathbb{Q}}
\newcommand{\R}{\mathbb{R}}
\newcommand{\C}{\mathbb{C}}
\newcommand{\F}{\mathbb{F}}
\newcommand{\A}{\mathbb{A}}
\newcommand{\PP}{\mathbb{P}}
\newcommand{\E}{\mathbb{E}}
\newcommand{\eps}{\varepsilon}
\newcommand{\dd}{\,\mathrm{d}}
\begin{document}
\section*{UnsolvedMath ID: 20001126}
\textbf{Source problem code:} catalogue entry
\subsection{Definitions and mathematical statement}
This source says ``Same problem, but over finite fields.'' The immediately preceding problem 3.13 defines a joint of lines in $\R^3$ as a point incident with three noncoplanar lines. Thus the recovered target is: given a prime power $q$ and $L$ distinct affine lines $\mathcal L$ in $\F_q^3$, determine the largest possible number of joints. A point is a joint if three incident lines have direction vectors spanning $\F_q^3$. Count points, not line triples or intersection multiplicities.

The supplied aid records a characteristic-independent bound of order $L^{3/2}$ and, attributing the sharper arbitrary-field result to Chao and Yu, the envelope
\[
 J(\mathcal L)\leq
 \min\left\{q^3,\ \frac{qL}{3},\
 \frac{L(\sqrt{8L+1}-3)}6\right\}.
\]
The first two terms retain the field-size and incidence ceilings. There is no field-size or characteristic restriction for this upper bound. Sharp $L^{3/2}$ order means that suitable configurations attain that order in a sufficiently large field; it is not a lower bound for every $L$ in a fixed small field.

The aid supplies two equality families. All axis-parallel lines have $L=3q^2$ and every one of the $q^3$ points is a joint, attaining the first two ceilings. For an integer $3\leq M\leq q$, take distinct $t\in\F_q$ and the $M$ planes
\[
 x+t y+t^2z=t^3.
\]
Their pairwise intersection lines number $L=\binom M2$, and their triple intersection points number $J=\binom M3$, attaining the third expression. This is an explicit interpretation of the aid's Vandermonde-plane family; the parameter restriction is essential.

The aid also reports a self-contained polynomial/deletion proof of
$J\leq(6^{1/3}+1)^{3/2}L^{3/2}$, addressing the positive-characteristic zero-derivative case by taking a $p$-th root over the perfect field $\F_q$. The quoted sharp theorem and the construction claims remain supplied results requiring independent verification; no characteristic-zero derivative argument is silently transferred to finite characteristic.
\subsection{Short English statement}
Count joints of distinct lines in finite three-dimensional space. The preceding real-space problem supplies the missing definition. The supplied aid gives a uniform sharp-type upper bound together with finite-field capacity ceilings and equality constructions in their stated parameter ranges.
\subsection{Sources}
\begin{itemize}
\item AIM list \emph{Recent trends in additive combinatorics}, problem 3.14 attributed to T. Tao: \url{https://aimath.org/WWN/additivecomb/additivecomb.pdf}.
\item Local registry \texttt{data/unsolvedmath/problems.json}, permanent ID 20001126, exact source code \texttt{catalogue entry}; canonical location \texttt{aim-combinatorics-notes.json notes[250]}. Companion ID 20001125, \texttt{AIM-COMBINATORICS-0250}, problem 3.13, \texttt{notes[249]}, supplies the real-space joints antecedent.
\item The attached aid attributes the exact arbitrary-field theorem to Chao and Yu, without full bibliographic details. Attempt 1, full solution, problem status at run solved, rigor proved, low novelty confidence; machine-generated, not peer reviewed, requiring independent verification before citation.
\end{itemize}
\end{document}
